\documentclass{article}

\usepackage[margin=1in]{geometry}
\usepackage{amsmath}
\usepackage{amssymb}
\usepackage{amsthm}
\usepackage[utf8]{inputenc}

\title{Finite Difference Methods for Options Pricing}
\author{Vahan Hovhannisyan}
\date{2026}

\begin{document}

\maketitle

\section{Introduction}
This tutorial covers finite difference methods for finance and options theory, specifically focusing on finding solutions to the Black-Scholes equation for a European put option. The underlying theory is based on \textit{Numerical Methods in Finance with C++} by M. Capi\'{n}ski and T. Zastawniak, but we will implement the solution in Python.

\section{Parabolic Partial Differential Equations}
First, we consider the general parabolic differential equation:
\begin{equation} \label{eq:parabolic}
\frac{\partial v(t,x)}{\partial t} = a(t,x) \frac{\partial^2 v(t,x)}{\partial x^2} + b(t,x)\frac{\partial v(t,x)}{\partial x} + c(t,x)v(t,x) + d(t,x)
\end{equation}
for $(t,x)\in [0,T]\times [x_l,x_u]$.

The problem includes the following boundary conditions:
\begin{itemize}
    \item \textbf{Terminal boundary condition:} $v(T,x) = f(x)$.
    \item \textbf{Lower boundary condition:} $v(t,x_l) = f_l(t)$.
    \item \textbf{Upper boundary condition:} $v(t,x_u) = f_u(t)$.
\end{itemize}

\section{The Black-Scholes Equation}
A typical example of a parabolic PDE is the Black-Scholes equation:
\begin{equation} \label{eq:black-scholes}
\frac{\partial u(t,z)}{\partial t} + \frac{\sigma^2}{2}z^2\frac{\partial^2 u(t,z)}{\partial z^2} + rz\frac{\partial u(t,z)}{\partial z} - r u(t,z) = 0
\end{equation}
Here, $u(t,S(t))$ represents a pricing function for an underlying asset with price $S(t)$ at time $t$.

This equation is a special case of the general parabolic PDE with the following parameters:
\begin{itemize}
    \item $a(t,z) = -\frac{\sigma^2}{2}z^2$
    \item $b(t,z) = -rz$
    \item $c(t,z) = r$
    \item $d(t,z) = 0$
\end{itemize}

For a European put option with an expiry date $T$ and a strike price $K$, the boundary conditions are:
\begin{itemize}
    \item $f(z) = u(T,z) = \max\{K-z,0\}$
    \item $fu(t) = u(t,z_u) = 0$
    \item $fl(t) = u(t,z_l) = e^{-r(T-t)}K$
\end{itemize}

To model this, you need data about the underlying asset (the risk-free interest rate and standard deviation) and the option itself (the strike price, time to expiration, and the lower/upper bounds for the spot price).

\section{The Explicit Scheme}



The explicit finite difference method operates on a finite set of points $\{(t_i, x_j):i=0,\dots,i_{\max}\text{ and } j=0,\dots,j_{\max}\}$.

The grid step sizes are defined as $\Delta t = \frac{T}{i_{\max}}$ and $\Delta x = \frac{x_u - x_l}{j_{\max}}$.

The explicit scheme relies on the following recursive formula:
\begin{equation}
v_{i-1,j}=A_{i,j}v_{i,j-1} + B_{i,j}v_{i,j} + C_{i,j}v_{i,j+1} + D_{i,j}
\end{equation}

The coefficients for this equation are given by:
\begin{itemize}
    \item $A_{i,j} = \frac{\Delta t}{\Delta x} \Big( \frac{b_{i,j}}{2}-\frac{a_{i,j}}{\Delta x}\Big)$
    \item $B_{i,j}=1-\Delta t c_{i,j} + \frac{2\Delta t a_{i,j}}{\Delta x^2}$
    \item $C_{i,j}=-\frac{\Delta t}{\Delta x}\Big( \frac{b_{i,j}}{2} + \frac{a_{i,j}}{\Delta x}\Big)$
    \item $D_{i,j}=-\Delta t d_{i,j}$
\end{itemize}

Because the values of $v_{i,j}$ can be determined from the boundary conditions, the computation starts from $i=i_{\max}$ and finishes at $i=0$. The boundary conditions establish that $v_{i_{\max},j} = f_j$, $v_{i-1,0} = f_{l,i-1}$, and $v_{i-1,j_{\max}} = f_{u,i-1}$.

\section{The Implicit Scheme}
Implicit schemes use different approximations for the partial gradients of $v$, meaning there is no explicit formula for $v_{i-1,j}$. Instead, the iterative formula is written in matrix form:
\begin{equation} \label{eq:implicit-scheme}
\mathbf{v}_{i-1} = \mathbf{B}_i^{-1} (\mathbf{A}_i\mathbf{v}_i + \mathbf{w}_i)
\end{equation}

The vectors and matrices are structured as follows:
\begin{itemize}
    \item $\mathbf{v}_i$ is a column vector containing the values $v_{i,1}$ through $v_{i,j_{\max}-1}$.
    \item $\mathbf{A}_i$ is a tridiagonal matrix with the coefficients $A_{i,j}$, $B_{i,j}$, and $C_{i,j}$ along its diagonals.
    \item $\mathbf{B}_i$ is a tridiagonal matrix with the coefficients $E_{i,j}$, $F_{i,j}$, and $G_{i,j}$ along its diagonals.
    \item $\mathbf{w}_i$ is a vector that incorporates the boundary conditions along with the $D_{i,j}$ terms.
\end{itemize}

The scalars used to construct these matrices for the implicit scheme are defined as:
\begin{itemize}
    \item $A_{i,j} = 0$, $B_{i,j} = 1$, $C_{i,j} = 0$
    \item $D_{i,j} = -\Delta t d_{i-1,j}$
    \item $E_{i,j}=-\frac{\Delta t}{\Delta x}\Big(\frac{b_{i-1,j}}{2}-\frac{a_{i-1,j}}{\Delta x}\Big)$
    \item $F_{i,j} = 1 + \Delta t c_{i-1,j} - \frac{2 \Delta t a_{i-1,j}}{\Delta x^2}$
    \item $G_{i,j} = \frac{\Delta t}{\Delta x}\Big(\frac{b_{i-1,j}}{2}+\frac{a_{i-1,j}}{\Delta x}\Big)$
\end{itemize}

Finally, a related algorithm called the Crank-Nicolson Scheme can be derived by overriding the $A$ through $F$ scalar coefficients defined here.

\end{document}